\documentclass[11pt]{amsart}

\usepackage[T1]{fontenc}
\usepackage{lmodern}
\usepackage{microtype}
\usepackage{amsmath,amssymb,amsthm}
\usepackage{booktabs}
\usepackage[hidelinks]{hyperref}

\hypersetup{
  pdftitle={The Harborth Constant of C9 plus C9}
}

\newtheorem{theorem}{Theorem}
\newtheorem{proposition}[theorem]{Proposition}
\newtheorem{lemma}[theorem]{Lemma}

\newcommand{\Z}{\mathbb Z}
\newcommand{\GL}{\operatorname{GL}}
\newcommand{\AGL}{\operatorname{AGL}}

\title[The Harborth constant of $C_9\oplus C_9$]
      {The Harborth Constant of $C_9\oplus C_9$}

\date{}

\subjclass[2020]{11B30, 20K01}
\keywords{Harborth constant, zero-sum subset, finite abelian group,
computer-assisted proof}

\begin{document}

\begin{abstract}
For a finite abelian group $G$, let $\mathsf g(G)$ be the least integer
$k$ such that every subset of $G$ with at least $k$ elements contains an
$\exp(G)$-element zero-sum subset. We prove
\[
  \mathsf g(C_9\oplus C_9)=17.
\]
An explicit $16$-element set gives the lower bound. The upper bound is an
exhaustive computer-assisted verification after affine normalization to
sets containing $\{(0,0),(1,0)\}$. The computation uses exact subsum
states, canonical prefixes under a stabilizer of order $108$, and exact
integer accounting of every searched or pruned completion. This confirms
the $n=9$ case of the Gao--Thangadurai conjecture.
\end{abstract}

\maketitle

\section{Introduction}

Let $C_n$ denote the cyclic group of order $n$. Motivated by Harborth's
lattice-point problem \cite{Harborth1973}, the Harborth constant of a
finite abelian group $G$ is
\[
  \mathsf g(G)=
  \min\left\{
    k:
    \begin{array}{l}
      \text{every }A\subseteq G\text{ with }|A|\geq k
      \text{ contains }B\subseteq A,\\
      |B|=\exp(G)\text{ and }\displaystyle\sum_{b\in B}b=0
    \end{array}
  \right\}.
\]
Gao and Thangadurai conjectured \cite{GaoThangadurai2004} that
\begin{equation}\label{eq:GT}
  \mathsf g(C_n\oplus C_n)=
  \begin{cases}
    2n-1,& n\text{ odd},\\
    2n+1,& n\text{ even}.
  \end{cases}
\end{equation}
The rank-two problem remains open in general; see
\cite[Section~1.1]{GuillotEtAl2019}. The parameter $n=9$ is the first odd
composite case, and to the best of our knowledge its exact value was not
previously determined.

\begin{theorem}\label{thm:main}
Every $17$-element subset of $C_9\oplus C_9$ contains a $9$-element
zero-sum subset, whereas some $16$-element subset contains none.
Consequently,
\[
  \mathsf g(C_9\oplus C_9)=17.
\]
\end{theorem}

The lower bound is elementary. The upper bound is the exact exhaustive
verification described below, which is recomputable from the data printed
in Section~\ref{sec:search}.

\section{Lower bound and affine normalization}

Throughout, write
\[
  R=\Z/9\Z,
  \qquad
  G=R^2,
  \qquad
  e_1=(1,0).
\]

\begin{proposition}\label{prop:lower}
One has $\mathsf g(G)\geq17$.
\end{proposition}

\begin{proof}
Let
\[
  A=\{0,1,\ldots,7\}\times\{0,1\}.
\]
The set $A$ has $16$ elements. Any $9$-element subset of $A$ contains
$t$ points with second coordinate $1$, where $1\leq t\leq8$, because
each horizontal layer has only eight points. Its second-coordinate sum is
therefore nonzero modulo $9$. Thus $A$ contains no $9$-element zero-sum
subset.
\end{proof}

Call $v\in G$ \emph{primitive} if $v\notin3G$.

\begin{lemma}\label{lem:primitive}
Every subset of $G$ with more than nine elements contains two points whose
difference is primitive.
\end{lemma}

\begin{proof}
The subgroup $3G$ has order $9$. If every pairwise difference in a set
$S$ belonged to $3G$, then, after fixing $s_0\in S$, one would have
\[
  S\subseteq s_0+3G,
\]
and hence $|S|\leq9$.
\end{proof}

\begin{lemma}\label{lem:normalization}
If a $17$-element subset of $G$ without a $9$-element zero-sum subset
exists, then one exists containing
\[
  B=\{0,e_1\}.
\]
\end{lemma}

\begin{proof}
By Lemma~\ref{lem:primitive}, such a set contains $a,b$ for which
$b-a$ is primitive. At least one coordinate of $b-a$ is a unit in $R$,
so there is an $M\in\GL_2(R)$ with
\[
  M(b-a)=e_1.
\]
The affine map $x\mapsto M(x-a)$ sends $a,b$ to $0,e_1$.

For any affine automorphism $\phi(x)=Mx+c$ and any nine points
$x_1,\ldots,x_9$,
\[
  \sum_{i=1}^9\phi(x_i)
  =
  M\sum_{i=1}^9x_i+9c
  =
  M\sum_{i=1}^9x_i.
\]
Since $M$ is invertible, a nine-term sum is zero if and only if its image
sum is zero.
\end{proof}

\section{Exhaustive verification}\label{sec:search}

Set $E=G\setminus B$ and order its $79$ elements increasingly by the
integer representative $x+9y$. Write
\[
  E=\{q_0,q_1,\ldots,q_{78}\}
\]
in this order. A normalized candidate is $B\cup T$ with
\[
  T\in\binom{E}{15}.
\]

\subsection{Exact subsum states}

For a partial selected set $P\subseteq G$, define
\[
  \Sigma_j(P)=
  \left\{
    \sum_{x\in U}x:
    U\subseteq P,\ |U|=j
  \right\},
  \qquad
  0\leq j\leq8,
\]
with $\Sigma_0(P)=\{0\}$.

\begin{lemma}\label{lem:update}
Assume that $P$ contains no $9$-element zero-sum subset. A point
$z\notin P$ may be added while preserving this property if and only if
\[
  -z\notin\Sigma_8(P).
\]
After adding $z$,
\[
  \Sigma_j(P\cup\{z\})
  =
  \Sigma_j(P)\cup
  \bigl(z+\Sigma_{j-1}(P)\bigr),
  \qquad
  1\leq j\leq8.
\]
\end{lemma}

\begin{proof}
Any newly created $9$-element zero-sum subset must contain $z$, and its
other eight points must sum to $-z$. The update formula partitions the
$j$-element subsets according to whether they contain $z$.
\end{proof}

The verifier represents each $\Sigma_j(P)$ by an exact $81$-bit set and
updates the states in descending order of $j$.

\subsection{Canonical prefixes}

Let
\[
  \Gamma=\operatorname{Stab}_{\AGL_2(R)}(B)
\]
be the setwise affine stabilizer of $B$. Its elements are exactly the maps
\[
  (x,y)\longmapsto(x+ay,uy)
  \quad\text{and}\quad
  (x,y)\longmapsto(1-x+ay,uy),
\]
where $a\in R$ and $u\in R^\times$. Hence
\[
  |\Gamma|=2\cdot9\cdot6=108.
\]

Let $\pi_\gamma$ be the permutation of the positions
$\{0,\ldots,78\}$ induced by $\gamma\in\Gamma$.

For $p=(p_1,p_2,p_3)$ with $p_1<p_2<p_3$, call $p$ \emph{canonical}
if, for every $\gamma\in\Gamma$, the increasing rearrangement of
\[
  \bigl(
    \pi_\gamma(p_1),
    \pi_\gamma(p_2),
    \pi_\gamma(p_3)
  \bigr)
\]
is lexicographically at least $p$.

\begin{lemma}\label{lem:canonical}
Every $\Gamma$-orbit in $\binom{E}{15}$ contains a set whose first three
positions form a canonical triple.
\end{lemma}

\begin{proof}
Choose the lexicographically least set
\[
  T=\{t_1<\cdots<t_{15}\}
\]
in its orbit. If $(t_1,t_2,t_3)$ were not canonical, some
$\gamma\in\Gamma$ would send these three positions, after sorting, to a
triple smaller than $(t_1,t_2,t_3)$. The first three positions of
$\gamma(T)$ are lexicographically no larger than that sorted triple,
contradicting the minimality of $T$.
\end{proof}

Exact enumeration gives $992$ canonical triples. A triple can be the first
three positions of a $15$-element set only if $p_3\leq66$, since twelve
larger positions must remain. Exactly $831$ canonical triples satisfy this
condition; these are the \emph{feasible canonical prefixes}. The exact
number of completions below them is
\begin{equation}\label{eq:expected}
  \sum_{\substack{p\text{ canonical}\\p_3\leq66}}
  \binom{78-p_3}{12}
  =
  811{,}509{,}569{,}804{,}714.
\end{equation}
The prefixes need not be unique orbit representatives;
Lemma~\ref{lem:canonical} is used only to ensure that no orbit is omitted.

\subsection{Search and verification}

At a search node, let $P$ be the current partial set and suppose that
$r$ further points are required. The verifier uses the following exact
rejection rules.

\begin{enumerate}
\item
A remaining point $z$ is individually inadmissible exactly when
\[
  -z\in\Sigma_8(P),
\]
by Lemma~\ref{lem:update}.

\item
If fewer than $r$ remaining points are individually admissible, no
completion exists.

\item
On the individually admissible points, form a graph in which distinct
$u,v$ are adjacent when
\[
  -(u+v)\notin\Sigma_7(P).
\]
Every valid $r$-point completion must induce an $r$-clique: a nonedge
$u,v$, together with seven points of $P$, would form a $9$-element
zero-sum subset. The verifier applies iterative degree peeling, which
preserves every $r$-clique, followed by an exact clique search, and
rejects a node only when no $r$-clique exists.
\end{enumerate}

The depth-first search partitions the lexicographic completion family at
each recursive node into disjoint subfamilies. When a subfamily is
rejected, its exact cardinality is credited using an integer binomial
coefficient. Thus every completion counted in \eqref{eq:expected} is
either examined or belongs to a soundly rejected subtree. The final
aggregate is as follows.

\begin{center}
\small
\begin{tabular}{@{}lr@{}}
\toprule
Quantity & Value \\
\midrule
Affine stabilizer size
  & $108$ \\
Canonical triples
  & $992$ \\
Feasible canonical prefixes
  & $831$ \\
Candidate completions
  & $811{,}509{,}569{,}804{,}714$ \\
Depth-first-search nodes
  & $879{,}672{,}298$ \\
Completions covered
  & $811{,}509{,}569{,}804{,}714$ \\
Valid completion found
  & no \\
\bottomrule
\end{tabular}
\end{center}

\begin{proposition}\label{prop:search}
There is no $T\in\binom{E}{15}$ for which $B\cup T$ contains no
$9$-element zero-sum subset.
\end{proposition}

\begin{proof}
The exact search rejected every completion below every feasible canonical
prefix, and the covered count equals the independently computed total in
\eqref{eq:expected}. By Lemma~\ref{lem:canonical}, every $\Gamma$-orbit
of normalized candidates has a searched representative.
\end{proof}

\paragraph{Exact verification.}
Section~\ref{sec:search} fixes the computation completely, so a referee can
recompute every number above from this paper alone. Ordering $E$ by $x+9y$
regenerates the $108$ maps of $\Gamma$ and the position permutations
$\pi_\gamma$, the $992$ canonical triples, the $831$ of them with
$p_3\leq66$, and the total \eqref{eq:expected}; the search is then driven by
Lemma~\ref{lem:update} and the three rejection rules, each rejected
subfamily being credited by an exact binomial coefficient, so that the
covered count is recomputed and compared with \eqref{eq:expected}. All
arithmetic is exact integer and bitset arithmetic, with no floating point
and no randomness, and the subsum states are self-tested against direct
enumeration before the search begins. The computation was re-run
independently on a separate machine, by a program written from these
definitions rather than from the original code, and the six quantities of
the table that do not depend on the traversal order agree with what is
reported here: the stabilizer size $108$, the $992$ canonical triples, the
$831$ feasible prefixes, the total \eqref{eq:expected}, its equality with
the completions covered, and the absence of a valid completion; only the
node count is specific to an implementation. The program and its transcript
are retained in an auxiliary archive, available on request; nothing above
depends on them. The re-run is a second computation, not a proof.

\begin{proof}[Proof of Theorem~\ref{thm:main}]
Proposition~\ref{prop:lower} gives the lower bound. If a $17$-element set
without a $9$-element zero-sum subset existed,
Lemma~\ref{lem:normalization} would produce one containing $B$,
contradicting Proposition~\ref{prop:search}. Hence
\[
  \mathsf g(C_9\oplus C_9)\leq17,
\]
and equality follows.
\end{proof}

\begin{thebibliography}{9}

\bibitem{GaoThangadurai2004}
W.~D. Gao and R. Thangadurai,
\emph{A variant of Kemnitz conjecture},
J. Combin. Theory Ser. A \textbf{107} (2004), no.~1, 69--86.

\bibitem{GuillotEtAl2019}
P. Guillot, L.~E. Marchan, O. Ordaz, W.~A. Schmid, and H. Zerdoum,
\emph{On the Harborth constant of $C_3\oplus C_{3p}$},
J. Th\'eor. Nombres Bordeaux \textbf{31} (2019), no.~3, 613--633.

\bibitem{Harborth1973}
H. Harborth,
\emph{Ein Extremalproblem f\"ur Gitterpunkte},
J. Reine Angew. Math. \textbf{262/263} (1973), 356--360.

\end{thebibliography}

\end{document}